\section{Introduction}
\label{sec:intro}

Skin cancer is among the most frequently diagnosed malignancies worldwide, and melanoma, though less common than the keratinocyte carcinomas, accounts for a disproportionate share of the associated mortality \cite{ref1}. Because prognosis depends strongly on the stage at detection --- lesions caught while still confined to the epidermis carry markedly more favourable outcomes than those found after invasion --- early and accurate identification of suspicious pigmented lesions is of considerable clinical importance. The diagnostic burden nonetheless falls on a specialist workforce that is small relative to the volume of lesions requiring assessment, which has drawn sustained attention to computational methods that support their examination.

Dermoscopy, a non-invasive technique that suppresses surface reflection to reveal subsurface pigmented structures, is now a routine modality in this assessment \cite{ref2}, and segmentation of the lesion region is a foundational step within it. The delineated region fixes the measured extent and shape of the lesion, which underpin the morphological criteria --- asymmetry, border irregularity, and diameter --- of established diagnostic rules \cite{ref3}, and it further conditions downstream classification systems that operate on the region or on features derived from it. Manual delineation, however, is laborious and varies appreciably between observers, most of all along the lesion margin, which motivates automated methods.

That margin is the crux of the task. Dermoscopic images frequently show low contrast between lesion and skin, so the transition is gradual rather than abrupt and cannot be located from local intensity gradients alone; margins are often diffuse or irregular, fading into adjacent skin over an extended transition zone. Structures that are not part of the lesion --- hair, air bubbles, ruler markings, and vascular patterns --- may traverse or obscure the boundary \cite{ref2}, and a single image may contain more than one pigmented structure, while lesions vary widely in scale, shape, colour, and texture. These conditions bear most heavily on the boundary, where the evidence separating lesion from background is weakest.

Automated approaches have developed along two lines. Convolutional encoder--decoder networks, archetypally the U-Net \cite{ref4}, join contracting and expanding paths by skip connections, and were later refined through attention gating of those connections \cite{ref5}, nested and densely connected decoders \cite{ref6}, and multi-scale aggregation; these capture local structure well, but the limited receptive field of convolution restricts their access to long-range context. Self-attention architectures answered this limitation, either replacing the convolutional encoder with a vision transformer \cite{ref7} or combining convolutional and transformer components in hybrid encoders \cite{ref8}, admitting global context at appreciably greater computational and data cost. More recently, large-scale visual foundation models --- most prominently the Segment Anything Model family \cite{ref9,ref10} --- have been adapted to medical segmentation, their image encoders supplying transferable representations that reduce the dependence on scarce annotated data.

Each line leaves difficulties that bear directly on boundary delineation. Foundation models, trained predominantly on natural images, are separated from dermoscopy by a domain gap \cite{ref11}, so that the adaptation mechanism rather than the backbone governs how faithfully fine structure is preserved. More fundamentally, most methods are both optimised and evaluated by region-overlap criteria. These are area-based, and the contour neighbourhood occupies only a small fraction of a compact lesion, so a contour displaced by a small distance alters the overlapping area only slightly \cite{ref12}; optimisation and evaluation on this basis exert little pressure on the very region where dermoscopic segmentation is hardest, and can fail to separate methods whose contours differ appreciably. Where boundary-sensitive content is treated at all, it is usually handled implicitly, merged with contextual information in a single undifferentiated representation dominated by the low-frequency interior.

Two remedies follow directly. Boundary-relevant content can be represented explicitly in the frequency domain, where the intensity transitions of lesion margins occupy the higher-frequency band \cite{ref13,ref14} and separate from the low-frequency interior more cleanly than in a single spatial representation; and it can be assessed explicitly with measures defined on the contour, since region overlap alone cannot characterise it. The Original MAFF-Net \cite{ref15} moved in this direction, adapting a frozen foundation backbone to dermoscopy through multi-scale perception adaptation and frequency-guided feature fusion with an auxiliary reconstruction branch, and demonstrating that spatial and spectral information can be integrated effectively. It does not, however, maintain the spatial and frequency representations as separate streams that are statistically reconciled before they are combined; closing that gap --- fusing the two domains as explicitly separated, reconciled streams --- motivates the present work.

Accordingly, this work investigates a targeted dual-stream refinement of the fusion stage of the MAFF-Net framework, rather than a replacement of the framework as a whole. The pretrained encoder, the convolutional decoder, and the auxiliary reconstruction branch are retained from the Original MAFF-Net, and the redesign is confined to the pathway in which spatial and spectral information are combined. The multi-scale features produced by the frozen encoder are routed simultaneously along a spatial stream, which preserves regional and contextual structure, and a frequency stream, which emphasises boundary-sensitive spectral content; the two are brought onto a common statistical footing and then merged by a learned gate whose weighting is inferred from the features themselves rather than fixed in advance. Training combines a region objective with auxiliary boundary-focused supervision applied to the segmentation outputs and to the reconstruction branch. The resulting model is evaluated against standard convolutional and transformer baselines under a protocol that reports region-based and boundary-based measures together, so that its effect on the lesion contour is assessed separately from its effect on the extent of the predicted region.

The contributions are threefold. First, region-overlap metrics alone are shown to be insufficient to characterise methods that act on the lesion margin: evaluating the models on the held-out test splits of three public dermoscopic benchmarks with region-based and boundary-based measures in combination yields different rankings under the two families of measures, so that boundary quality must be assessed with contour-based measures rather than inferred from region overlap. Second, the proposed dual-stream model --- a parameter-free Frequency Branch, per-stream Feature Alignment, and a learned Gated Fusion on a frozen foundation backbone --- attains higher region-based and boundary-based scores than four standard convolutional and transformer baselines on all three benchmarks under an identical single-run protocol. Third, the dual-stream reformulation of the fusion pathway is investigated, and a leave-one-out ablation isolates the marginal role of its three components: disabling any one degrades all boundary measures while leaving region overlap essentially unchanged on the dataset examined, confining their effect to boundary fidelity rather than to the extent of the predicted region. The encoder, decoder, and reconstruction branch are adopted from the Original MAFF-Net and are not claimed as contributions of this work; the boundary-centric supervision under which the network is trained is described in Section III-H as part of the training objective.

The remainder of the paper is organised as follows. Section II reviews related work on dermoscopic lesion segmentation; Section III presents the proposed methodology; Section IV specifies the datasets, implementation, evaluation measures, and experimental protocol; Section V reports the results; Section VI discusses their interpretation and the limitations of the study; and Section VII concludes.
